\section*{Chapitre \ref{ch:fm:technology-strategy} --- Stratégie technologique}

\begin{solution}{exo:fm:technology-strategy:1}
19~\%, de 209.4~millions de dollars à 249.2~millions de dollars. En part du chiffre d'affaires~: 13.8~\%, 6.6~\%, 7.5~\%, 9.3~\%, 10.1~\%, 8.2~\% et 6.9~\% de 2019 à 2025~: la
part varie avec le chiffre d'affaires parce que le coût, lui, ne varie pas.
\end{solution}

\begin{solution}{exo:fm:technology-strategy:2}
Seule~: $60(0.8+0.2/5)=50.4$ millions de dollars par an. Avec deux autres~: $60(0.8/3+0.2/5)=18.4$ millions de dollars.
\end{solution}

\begin{solution}{exo:fm:technology-strategy:3}
36\,000~\$, 450\,000~\$ et 946\,000~\$ par an~: 12.6~\%, 33.3~\% et 21.3~\% du total des coûts des paliers.
\end{solution}

\begin{solution}{exo:fm:technology-strategy:4}
Le market making sur contrats à terme $\lfloor48/4.16\rfloor=11$, plafonné à 5~; l'arbitrage sur ETF $\lfloor27/4.16\rfloor=6$, plafonné à 5~; le market making sur options 4~;
le crypto trade 2 au palier 2 ($\lfloor3/1.064\rfloor$)~; arbitrage statistique 1 au palier 2 ($\lfloor2/1.064\rfloor$). La simulation le confirme.
\end{solution}

\begin{solution}{exo:fm:technology-strategy:5}
Seuls 10~\% de son chiffre d'affaires se jouent dans les courses~: les remporter toutes ajoute au plus 2~millions de dollars par an, moins que les 4.16~millions de dollars que coûte le passage
au palier 3. Il gagne 3.31~millions de dollars au palier 1 et 1.15~millions de dollars au palier 3.
\end{solution}

\begin{solution}{exo:fm:technology-strategy:6}
Renouveler le contrat de baie, remplacer le commutateur en fin de vie et la facture de données de marché relèvent du run~; la nouvelle passerelle vers la plateforme est du change.
\end{solution}

\begin{solution}{exo:fm:technology-strategy:7}
Les cinq sociétés continuent toutes à se disputer le market making sur contrats à terme~; l'arbitrage sur ETF tombe à trois concurrents et le market making sur options à deux. Le total des dépenses de vitesse
passe de 68.6~millions de dollars à 86.9~millions de dollars et la part gaspillée de 89.6~\% à 91.8~\%~: une vitesse plus chère ne met pas fin à la course là où le pot est grand.
\end{solution}

\begin{solution}{exo:fm:technology-strategy:8}
Le gaspillage relève de l'ensemble du secteur, pas d'une seule société~: une société qui cesse de se mettre à niveau dans une stratégie où les autres se disputent la course perd le revenu de course ($\alpha V/m$) et
ne garde que sa part du reste. La bonne décision se prend stratégie par stratégie~: il faut concourir là où $\alpha V/m$ dépasse le saut de coût, et renoncer là où il ne le dépasse pas.
\end{solution}

\begin{solution}{pb:fm:technology-strategy:1}
\begin{enumerate}
\item Voir les \cref{def:fm:technology-strategy:tier,def:fm:technology-strategy:race}.
\item Un coût qui a augmenté de 19~\% en sept ans alors que le chiffre d'affaires a plus que doublé, puis a été divisé par deux~: un engagement fixe.
\item Environ une course par minute et par symbole, 5--10 millionièmes de seconde, un cinquième du volume, six sociétés concentrant plus de 80~\% des victoires et des défaites, un tiers du
coût de la liquidité.
\item 0.29, 1.35 et 4.45~millions de dollars par an et par stratégie~; les frais de plateforme viennent des grilles publiées dans le livre~14, le reste relève des hypothèses.
\item $V(\alpha\,\mathbf 1\{\text{premier}\}/m+(1-\alpha)/n)-c$.
\item Voir la \cref{fig:fm:technology-strategy:profits}~: palier 3 pour le market making sur contrats à terme, l'arbitrage sur ETF et le market making sur options~; palier 2 pour le
crypto trade~; palier 1 pour l'arbitrage statistique.
\item Aucun rival n'est plus rapide que le palier 1, donc le palier 2 remporte toutes les courses pour 1.06~millions de dollars de plus~; le palier 3 n'ajoute rien.
\item Une part de course plus élevée, un pot plus grand ou une vitesse moins chère.
\item Voir la \cref{prop:fm:technology-strategy:count}.
\item Cinq au palier 3 en market making sur contrats à terme et en arbitrage sur ETF (22.2~millions de dollars chacun)~; un au palier 2 en arbitrage statistique (2.5~millions de dollars)~; quatre au
palier 3 en market making sur options (18.1~millions de dollars)~; deux au palier 2 sur le crypto trade (3.6~millions de dollars).
\item 93.6~\%, 93.6~\%, 42.7~\%, 92.1~\% et 59.8~\%~; 89.6~\% au total.
\item Une société qui s'arrête perd son revenu de course au profit de celles qui continuent.
\item Les marchés continus créent des rentes d'arbitrage que la concurrence ne supprime pas~: elle ne fait qu'élever la vitesse nécessaire pour les capter, dans une course socialement
gaspilleuse~; ils proposent des enchères par lots fréquentes.
\item Voir la \cref{def:fm:technology-strategy:teams} et la \cref{fig:fm:technology-strategy:teams}.
\item Voir la \cref{def:fm:technology-strategy:runchange}.
\item 16.2~millions de dollars de run, 4.0~millions de dollars de change~: 80.2~\% de run.
\item Des charges technologiques de 70.6~millions d'euros en 2025, contre 66.6~millions d'euros en 2024, face à un produit net de trading de 485.8~millions d'euros.
\item Chaque année, sur les paliers des concurrents, la part de course mesurée à partir des exécutions propres de la société et des prix des paliers.
\item Palier 3 pour le market making sur contrats à terme, l'arbitrage sur ETF et le market making sur options, palier 2 pour le basis trade crypto, palier 1 pour
l'arbitrage statistique~; 89.6~\% des dépenses de vitesse du secteur sont du gaspillage par rapport à un scénario où personne ne se mettrait à niveau.
\item Ne courir que là où les courses décident de l'essentiel d'un grand pot, sur une plateforme partagée à faible latence~; revoir chaque année le palier de chaque stratégie et
intégrer chaque palier comme un coût de run auquel la société s'engage.
\end{enumerate}
\end{solution}

\begin{solution}{iq:fm:technology-strategy:1}
Une \omterm{def:fm:technology-strategy:teams}{équipe produit} construit et fait tourner le système de trading d'une activité et se mesure à ses résultats~; une \omterm{def:fm:technology-strategy:teams}{équipe plateforme} fournit des services partagés tels que
les données de marché et les passerelles, et se mesure à leur fiabilité et à la livraison des \omterm{def:fm:technology-strategy:teams}{équipes produit}.
\iqlookfor{responsabilité et indicateurs.}
\end{solution}

\begin{solution}{iq:fm:technology-strategy:2}
Lorsque les courses qu'elle remporte valent moins que le surcoût de la vitesse, ou lorsque les concurrents sont de toute façon plus rapides et qu'elle ne remporte aucune course.
\iqlookfor{part de course multipliée par le pot, comparée au coût.}
\end{solution}

\begin{solution}{iq:fm:technology-strategy:3}
$\alpha V=27$~; chaque société rapide touche $27/m$~; $27/m\ge4$ pour $m\le6$, donc les cinq.
\iqlookfor{l'argument de comptage.}
\end{solution}

\begin{solution}{iq:fm:technology-strategy:4}
Mesurer la part du revenu de la stratégie perdue face aux concurrents plus rapides dans les courses, estimer ce que le matériel permettrait de récupérer, et comparer avec son coût complet, ingénieurs compris~; décider
sur cette base, et non sur la seule latence.
\iqlookfor{revenu en jeu comparé au coût complet.}
\end{solution}

\begin{solution}{iq:fm:technology-strategy:5}
Chaque projet de change crée une infrastructure, des licences et un support qu'il faut faire tourner tant que l'activité existe, si bien que le change devient du run.
\iqlookfor{les engagements s'accumulent.}
\end{solution}

\begin{solution}{iq:fm:technology-strategy:6}
Elles regrouperaient les ordres à intervalles courts pour que de minimes écarts de vitesse ne décident plus de qui négocie~: $\alpha$ tend vers zéro, moins de sociétés
se disputent la course, et les dépenses de vitesse baissent.
\iqlookfor{la part de course comme levier.}
\end{solution}
